% ============================================================
% appendix/app-glossary.tex —— 附录 D 术语中英对照表
% 以 CONTRACT §11 为准，按概念分组，共 57 行。
% ============================================================

\section{术语中英对照表}
\label{app:glossary}

表~\ref{app:tab:glossary} 按“基础—相关—波函数—DFT—半经验/力场—性质”的顺序排列，共 57 条。中文词以 CONTRACT §11 为准；同一概念在全书只用这里给出的译名，首次出现时写成“中文（English，缩写）”。

{\small
\begin{longtable}{@{}p{3.4cm} p{11.2cm}@{}}
\caption{术语中英对照表（57 条）。}\label{app:tab:glossary}\\
\toprule
中文 & English \\
\midrule
\endfirsthead
\multicolumn{2}{@{}l}{\small 表~\ref{app:tab:glossary}（续）}\\
\toprule
中文 & English \\
\midrule
\endhead
\bottomrule
\endlastfoot
波函数 & wavefunction \\
多体薛定谔方程 & many-body Schrödinger equation \\
玻恩--奥本海默近似 & Born--Oppenheimer approximation \\
斯莱特行列式 & Slater determinant \\
基组 & basis set \\
自洽场 & self-consistent field (SCF) \\
平均场 & mean field \\
电子相关 & electron correlation \\
动态相关 & dynamic correlation \\
静态相关 & static (nondynamic) correlation \\
尺寸一致性 & size consistency \\
尺寸广延性 & size extensivity \\
变分性 & variationality \\
自旋污染 & spin contamination \\
微扰理论 & perturbation theory \\
组态相互作用 & configuration interaction (CI) \\
耦合簇 & coupled cluster (CC) \\
活性空间 & active space \\
自然轨道 & natural orbital \\
占据数 & occupation number \\
冻结核 & frozen core \\
交换相关 & exchange-correlation \\
泛函 & functional \\
自相互作用误差 & self-interaction error (SIE) \\
范围分离 & range separation \\
双杂化 & double hybrid \\
紧束缚 & tight-binding \\
力场 & force field \\
色散校正 & dispersion correction \\
色散 & dispersion \\
辅助基组 & auxiliary basis set \\
基组重叠误差 & basis set superposition error (BSSE) \\
势能面 & potential energy surface (PES) \\
几何优化 & geometry optimization \\
构象 & conformer \\
极化 & polarization \\
Ewald 求和 & Ewald summation \\
周期性边界条件 & periodic boundary conditions (PBC) \\
热化学 & thermochemistry \\
反应能 & reaction energy \\
活化能 & activation energy \\
电子亲和能 & electron affinity (EA) \\
电离能 & ionization energy (IE) \\
带隙 & band gap \\
梯度 & gradient \\
拉普拉斯 & Laplacian \\
动能密度 & kinetic energy density \\
等轨道指示器 & iso-orbital indicator \\
绝热连接 & adiabatic connection \\
交换相关穴 & exchange-correlation hole \\
氢键 & hydrogen bond \\
卤键 & halogen bond \\
溶剂化 & solvation \\
隐式溶剂模型 & implicit solvation model \\
雅各布天梯 & Jacob's ladder \\
复合方法 & composite method \\
决策树 & decision tree \\
\end{longtable}
}

\noindent{\footnotesize 说明：缩写只在正文首次出现时给出，之后直接用缩写；例如“密度泛函理论（density-functional theory，DFT）”之后写 DFT。表中“尺寸一致性”与“尺寸广延性”是不同概念：前者要求超体系能量等于碎片之和，后者要求能量随粒子数线性标度；截断 CI 两者都不满足，CC 在截断下仍保持广延性。}
